\documentclass[12pt,a4paper]{article}
\usepackage{graphicx} % Required for inserting images
\setlength{\topmargin}{0.0in}
\setlength{\oddsidemargin}{0.33in}
\setlength{\textheight}{9.0in}
\setlength{\textwidth}{6.0in}
\renewcommand{\baselinestretch}{1.25}


\usepackage{hyperref}
% \usepackage[utf8]{inputenc}
% \usepackage[T1]{fontenc}
\usepackage{float}
\usepackage{enumitem}
\usepackage{amsmath}
\usepackage{amssymb}
\usepackage{amsthm}
\usepackage{tikz}
\usepackage{amsfonts}
\usepackage{graphicx}
\usepackage{subcaption}
\usepackage{caption}
\newcommand{\dv}[2]{\frac{\mathrm{d} #1}{\mathrm{d} #2}}
\usepackage{caption}
\captionsetup[figure]{font=small} 
% \captionsetup[table]{font=small}
\usepackage[capitalize,nameinlink]{cleveref}
\usepackage{mathtools}
\newcommand{\defeq}{\vcentcolon=}
\newcommand{\eqdef}{=\vcentcolon}



\newtheoremstyle{boldthm}
  {3pt}        % space above
  {3pt}        % space below
  {\itshape}   % body font
  {}           % indent
  {\bfseries}  % theorem head font
  {.}          % punctuation
  {.5em}       % space after theorem head
  {\thmname{#1}\thmnumber{ #2}\thmnote{ \textbf{(#3)}}}

\newtheoremstyle{bolddefn}
  {3pt}
  {3pt}
  {\itshape} 
  {}
  {\bfseries}
  {.}
  {.5em}
  {\thmname{#1}\thmnumber{ #2}\thmnote{ \textbf{(#3)}}}


\theoremstyle{boldthm}
\newtheorem{theorem}{Theorem}[section]
\newtheorem{corollary}[theorem]{Corollary}
\newtheorem{lemma}[theorem]{Lemma}
\newtheorem{proposition}[theorem]{Proposition}

\theoremstyle{bolddefn}
\newtheorem{definition}[theorem]{Definition}
\newtheorem{example}[theorem]{Example}
\newtheorem{remark}[theorem]{Remark}

% Define environments
% \newtheorem{theorem}{Theorem}
% \newtheorem{definition}{Definition}
% \newtheorem{lemma}{Lemma}
% \newtheorem{corollary}{Corollary}


\title{Our Functional-Analytic Setting!}
\author{}
\date{}

\begin{document}

\maketitle


\section{Matt's Setting and the Issue it Poses for Us}

Throughout his work, Matt has identified $H$ with its dual. Thus, he considers $A : D(A) \subset H \rightarrow H$ and writes
\begin{align}
    &\operatorname{Sp}(A) = \{z \in \mathbb{C} : (A-zI) \text{ is not boundedly invertible}\},\\
    &\sigma_{\inf}(A) = \inf \{ \|Ax\| : x \in D(A), \|x\| = 1 \}, \\
    &\gamma(z, A) = \|(A-zI)^{-1} \|^{-1} = \min \{ \sigma_{\inf}(A-zI), \sigma_{\inf}(A^* - \overline{z}I)\},\\
    &\gamma(z, A) \le \operatorname{dist}(z, \operatorname{Sp}(A)),\\
     &\operatorname{Sp}_{\epsilon}(A) = \{z \in \mathbb{C} : \gamma(z, A) = 0\},\\
    &\operatorname{Sp}_{\epsilon}(A) = \{z \in \mathbb{C} : \gamma(z, A) < \epsilon\},\\
    & (A-zI)^*(A-zI), \text{ and } (A-zI)(A-zI)^*.
\end{align}

 In our work, we consider operators that map from a Hilbert space to its antidual, $A: H \rightarrow H^x.$ As a result, we cannot write $A-zI$ without properly specifying what $I$ is. Thus, we need to re-think what $A-zI$ is and what the correct notion of spectra and pseudospectra is for our work. 


\section{The Antidual Adjoint}

We first need to define the antidual adjoint. If $f \in H^x$ then $f$ is a bounded conjugate-linear functional, i.e., $f(au) = \overline{a} f(u).$ We want to define a canonical map $J_{\times} : H \rightarrow H^{\times}$ so that $J_{\times} u \in H^{\times \times}.$ To do so, we write 
\begin{equation}
    (J_{\times} u)(f) = \overline{f(u)}.
\end{equation}
Then, $(J_{\times}u)(af) = \overline{a f(u)} = \overline{a} (J_{\times} u)(f).$

Consider $A : H \rightarrow H^{\times}.$ We define the Banach adjoint $A^B : H^{\times \times} \rightarrow H^{\times}$ by 
\begin{equation*}
    \langle A^B f, u \rangle_{H^{\times}, H} = \langle f , Au \rangle_{H^{\times \times}, H^{\times}}.
    \end{equation*}
Using the canonical map $J_{\times},$ we can define the antidual adjoint $A^x: H \rightarrow H^{\times}$ via $A^{\times} = A^B J_{\times}.$ In other words, the antidual adjoint satisfies 
\begin{equation}
    \langle A^{\times} v, u \rangle_{H^{\times}, H} = \overline{ \langle Au, v \rangle_{H^{\times}, H}}.
\end{equation}
Going forward, we will drop the subscripts on the duality action where possible, for ease of notation. 


\section{Choosing the `identity' $M$}

Instead of forming $A-zI$ we will form $A-zM$ where
\begin{equation}
    M : H \rightarrow H^{\times} \qquad \langle Mu, v \rangle_{H^{\times}, H} \defeq (u, v)_{L^2},
\end{equation}
which is both positive definite and self-adjoint with respect to the antidual adjoint, i.e., $\langle Mu, u \rangle > 0$ for $u \ne 0$ and $M^{\times} = M.$ As a result, $M$ defines a metric.


When it exists, $M^{-1}$ defines a metric via the following inner product:
\begin{equation}
    M^{-1}: H^{\times} \rightarrow H, \qquad (f, g)_{M^{-1}} \defeq \langle f, M^{-1} g \rangle_{H^{\times}, H}.
\end{equation}

See Homework:
\begin{itemize}
    \item $M$ is self-adjoint,
    \item $M^{-1}$ and $M$ induce metrics,
    \item New notion of orthogonality using these metrics. 
\end{itemize}


\section{Injection Modulus w.r.t $M$ and $M^{-1}$}

With this $M,$ we can consider $A-zM$ instead of $A-zI$ for spectral and pseudospectra computations. This choice will be clarified below. 

\subsection{A New Injection Modulus}
The injection modulus measures how injective an operator is. It can be thought of as the infinite-dimensional analog of a singular value. Matt defines it as follows:
\begin{equation}
    \sigma_{\inf}(A) = \inf \{ \|Ax\| : x \in D(A), \|x\| = 1 \}.
\end{equation}

We want to measure the injection modulus of $B_z \defeq A - zM : H \rightarrow H^{\times}.$ To do so, we will use the metrics induced by $M^{-1}$ and $M.$ This choice yields the following:
\begin{equation}
    \sigma_{\inf}^2(B_z) = \inf_{u \ne 0} \frac{ \|B_z u \|_{M^{-1}}^2}{\|u\|_{M}^2}.
\end{equation}
Expanding the numerator out, we obtain
\begin{equation}
    \langle B_z u, M^{-1} B_z u \rangle_{H^{\times}, H}.
\end{equation}
Using the antidual adjoint, $\langle B_z u, v \rangle_{H^{\times}, H} = \overline{ \langle B_z^{\times} v, u \rangle_{H^{\times}, H}},$ we can write
\begin{equation}
    \langle B_z u, M^{-1} B_z u \rangle_{H^{\times}, H} = \overline {\langle B_z^{\times} (M^{-1} B_z u), u \rangle_{H^{\times}, H}} .
\end{equation}
Notice that $\langle B_z u, M^{-1} B_z u \rangle_{H^{\times}, H} = \| B_xu\|_{M^-1}^2$ is real and nonnegative. Thus, we can drop the conjugate and write
\begin{equation}
    \sigma_{\inf}^2(B_z) = \inf_{u \ne 0} \frac{{\langle B_z^{\times} (M^{-1} B_z u), u \rangle_{H^{\times}, H}}}{{\langle Mu, u \rangle_{H^{\times}, H}}},
\end{equation}
or
\begin{equation}
    \sigma_{\inf}^2(A-zM) = \inf_{u \ne 0} \frac{{\langle (A-zM)^{\times} M^{-1} (A-zM) u, u \rangle_{H^{\times}, H}}}{{\langle Mu, u \rangle_{H^{\times}, H}}}.
\end{equation}

\section{Connection to Matt's Pseudospectra (Issues)}

For our work, we define $M$ by $\langle Mu, v \rangle_{H^{\times}, H} = (u, v)_{L^2}.$ To get Matt's operator $\mathcal{A} : H \rightarrow H$ that relates to our operator $A : H \rightarrow H^{\times},$ we simply need to apply $M^{-1}$ to $A:$
\begin{equation}
    \mathcal{A} = M^{-1} A : H \rightarrow H.
\end{equation}
\textbf{QUESTIONS}:
\begin{itemize}
    \item Does it matter that Matt's operator is from $D(\mathcal{A}) \subset H$ to $H$???
    \item Does it matter that we use $L^2$ and Matt uses $H$? I think it works if we define $M$ to give the $H$ inner product.
    \item $M$ is a mass operator not a Riesz map. Does $M^{-1}$ exist? $M$ is not onto for $H^{\times}$? Can we say things like $A = M \mathcal{A}$? I am hesitant to treat $M^{-1}$ like a Riesz map.

\end{itemize}
In this way, $M^{-1}(A-zM) = \mathcal{A} - zI$ where $I : H \rightarrow H.$ Furthermore, notice that with our choice of $M$ we have 
\begin{equation}
    \| M w \|_{M^{-1}}^2 = \langle Mw, w \rangle = \|w \|_M^2 = \|w \|_{L^2}^2.
\end{equation}
Thus,
\begin{equation}
    \|(A-zM) u \|_{M^{-1}}^2 = \|M(\mathcal{A} - zI) u \|_{M^{-1}}^2 = \| \mathcal{A} - zI\|_{L^2}^2.
\end{equation}
Hence,
\begin{equation}
    \sigma_{\inf}(A-zM) = \sigma_{\inf}(\mathcal{A}-zI).
\end{equation}

This also holds for the adjoint injection modulus. The Hilbert adjoint of $\mathcal{A}$ satisfies 
\begin{equation}
    (\mathcal{A}u, v)_{L^2} = (u, \mathcal{A}'v)_{L^2}.
\end{equation}
Notice that
\begin{equation}
    (\mathcal{A}u, v)_{L^2} = \langle M(\mathcal{A} u), v \rangle = \langle M(M^{-1} Au), v \rangle = \langle Au, v \rangle.
\end{equation}
Using the antidual adjoint, we have 
\begin{equation}
    (\mathcal{A}u, v)_{L^2} = \overline{ \langle A^{\times} v, u \rangle}
\end{equation}
Write $w = M^{-1} A^{\times} v \in H.$ Then $Mw = A^{\times} v$ and 
$\overline{ \langle Mu, v \rangle} = \overline{(w, u)_{L^2}} = (u, w)_{L^2}.$ Thus,
\begin{equation}
    \langle Au, v \rangle = \overline {\langle Mw, u \rangle} = (u, w)_{L^2} = (u, M^{-1} A^{\times} v)_{L^2}. 
\end{equation}
Thus, $\mathcal{A}' = M^{-1} A^{\times}.$

Furthermore, 
$\mathcal{A}' - \overline{z} I $ = $M^{-1}(A^{\times} - \overline{z}M).$ Hence,
\begin{align}
    \sigma_{\inf}(\mathcal{A}' - \overline{z}I) &= \inf_{\|u\|_{L^2} = 1} \| M^{-1} (A - zM)^{\times} u \|_{L^2} \\
    & = \inf_{\|u\|_{L^2} = 1} \| (A - zM)^{\times} u \|_{M^{-1}}\\
    &= \sigma_{\inf}(A-zM)
 \end{align}

\textbf{ AM I ACCIDENtALLY TREATING $M$ AS A RIESZ MAP !!??}

\subsection{Poisson}

\begin{itemize}
    \item Take $Lu = -u''$ and $D(L) = H^2(0, \pi) \cap H_0^1(0, \pi)$ which acts like (?) $L : D(L) \subset L^2(0, \pi) \rightarrow L^2(0, \pi).$
    \item Write $A : H_0^1(0, \pi) \rightarrow (H_0^1(0, \pi))^{\times} \defeq H^{-1}(0,\pi),$ defined by $\langle Au, v \rangle = (u', v')_{L^2}.$
    \item Define the mass map $M : H_0^1(0, \pi) \rightarrow H^{-1}(0, \pi)$ by $\langle Mu, v\rangle = (u, v)_{L^2}.$

    \item The `weak' shift operator is $A-zM$. For sufficiently regular $u,$ 
\begin{equation}
    \langle Au, v \rangle = (Lu, v)_{L^2},
\end{equation}
via integration by parts. Also, 
\begin{equation}
    \langle z Mu, v \rangle = (zu, v)_{L^2}.
\end{equation}
Thus, 
\begin{equation}
    \langle (A-zM)u, v\rangle = ((L-zI)u, v)_{L^2}.
\end{equation}
In other words,
\begin{equation}
    (A-zM)u = M((L-zI)u).
\end{equation}

\item \textbf{THIS IS WHERE THE ISSUE IS!!} We need $(L-zI)u \in H_0^1(0, \pi).$ Should we instead define another map 
$$J : L^2 \rightarrow H^{-1}(0, \pi)$$ by 
$$\langle Jf, v \rangle = (f, v)_{L^2}$$
so that $M = J |_{H_0^1(0, \pi)}$? Then $J^{-1}(A-zM)u = (L-zI)u$????

% \item Now, 
% \begin{align}
%     \|(A-zM)u\|_{M^{-1}}^2  &= \langle (A-zM) u, M^{-1} (A-zM) u \rangle \\
%     &= \langle M((L-zI)u), (L-zI)u \rangle \\
%     &= \| (L-zI) u \|_{L^2}^2.
% \end{align}

% \item Thus,
% \begin{equation}
%     \sigma_{\inf}(L-zI) = \inf_{u \ne 0} \frac{\|(A-zM)u \|_{M^{-1}}}{\|u\|_{M}}.
% \end{equation}

\item Now, 
\begin{align}
    \|J^{-1}(A-zM)u\|_{L^2}^2 
    &= \| (L-zI) u \|_{L^2}^2  ???
\end{align}

\item Thus,
\begin{equation}
    \sigma_{\inf}(L-zI) = \inf_{u \ne 0} \frac{\|(J^{-1} A-zM)u \|_{L^2}}{\|u\|_{L^2}} ???
\end{equation}

\end{itemize}


\section{Discretization}

We discretize $B_z = A-zM$ with $B_h.$ We say that $B_h$ acts between $(\mathbb{C}, \| \cdot \|_{M_h})$ and $(\mathbb{C}, \| \cdot \|_{M_h^{-1}})$
where $\|x\|_{M_h}^2 = x^*M_h x$ and $\|y\|_{M_h^{-1}}^2 = y^*M_h^{-1} y.$

In this way, we build the generalized SVD of $B_h,$
\begin{equation}
    B_h = U \Sigma (V^* M_h),
\end{equation}
where $U^* M_h^{-1} U = I,$ $V^* M_h V = I,$ and $\Sigma$ holds the singular values with respect to $M_h$ and$ M_h^{-1}$, i.e.,
\begin{equation}
    \sigma_j(B_h)  = \sqrt{\lambda_j(B_h^* M_h^{-1} B_h, \, M_h)}.
\end{equation}
From this, we can show that 
\begin{equation}
    \| B_h^{-1} \|_{M^{-1}, M} = 1/\sigma_{\min}(B_h),
\end{equation}
 where $\|B_h\|_{M^{-1}, M} = \sup_{f \ne 0} \| B_h^{-1} f \|_{M_h^{-1}} / \|f \|_{M_h}.$


 \noindent
See Homework:
\begin{itemize}
    \item SVD proof,
    \item Smallest singular value proof,
    \item Orthogonal/Isometry proof.
\end{itemize}

\end{document}